\chapter{関連内容}

ここには，関連内容を書きます．
来年の方に，引き継ぐなど，より提案手法がわかりやすくなるように
手法で使用する技術の説明などを行います．

たとえば，群知能のアントコロニー最適化などに関する説明などです．
そして，第\ref{sec:chap_4}章で自分の手法を提案します．


\section{アントコロニー最適化}


\subsection{群知能}

画像は以下のように入れることができます．
ラベルをつけると，他の場所\ref{fig:cute}から番号で参照でき，便利です．
（わざわざあとですべての図の番号を自分で変えなくて良いので）

tbpと指定するのが標準・学会デファクトらしいです．

\begin{figure}[tbp]
    \begin{center}
        \includegraphics[width=0.8\linewidth]{images/chapter3/cute.jpeg}
    \end{center}
    \caption{cute girl}
    \label{fig:cute}
\end{figure}

\section{LaTeX}

LaTeXでは，数式を利用することができます．

本文中なら\(x\)，中央寄せなら\[y = x^2\]のようにすると良いです．

また，数式に番号をつけたい場合は\ref{eq:x3}のようにすると良いです．

\begin{equation}
    y = x^3
    \label{eq:x3}
\end{equation}

\section{箇条書き}

箇条書きも利用できます．

\begin{itemize}
    \item this
    \begin{itemize}
        \item 入れ子
        \item もできます
    \end{itemize}
    \item is
\end{itemize}

数字付きも利用できます．

\begin{enumerate}
    \item this is 1
    \item thsi is 2
\end{enumerate}